%% ASPLOS 2026 submission template
%% Class:  acmart (sigplan proceedings)
%% Final project report.

\documentclass[sigplan,nonacm]{acmart}

\settopmatter{printfolios=true,printccs=false,printacmref=false}

%% ---- Packages -----------------------------------------------------------
\usepackage{booktabs}
\usepackage{subcaption}
\usepackage{balance}
\usepackage{listings}
\usepackage{xurl}

%% Breakable code spans for long file paths and identifiers.
%% \texttt{...} cannot break at _, /, ::; \path{...} can.
\newcommand{\code}[1]{\texttt{\nolinkurl{#1}}}

\emergencystretch=5em
\hbadness=10000
\tolerance=2000

\lstset{
  basicstyle=\ttfamily\footnotesize,
  breaklines=true,
  columns=fullflexible,
  frame=none,
  xleftmargin=1.5em,
  numbers=none
}

%% ---- Metadata ------------------------------------------------------------
\acmConference[ASPLOS '26]{31st International Conference on Architectural
  Support for Programming Languages and Operating Systems}{March 22--26,
  2026}{Pittsburgh, PA, USA}

%% ---- Title ---------------------------------------------------------------
\title{Optimizing Qwen2.5-Coder Decode on Tenstorrent Blackhole:\\
       Two Wins, One Dead End, and What We Learned}

\author{Rassul Magauin}
\author{Assylzhan Khamiyev}
\author{Mohammed Rashed Ali Yahmoor Alshehhi}
\author{Sultan Akimaliyev}
\author{Hamad Khalifa Alyahyaee}
\affiliation{%
  \institution{Mohamed bin Zayed University of Artificial Intelligence}
  \city{Abu Dhabi}\country{UAE}}

\renewcommand{\shortauthors}{Magauin et al.}

%% ---- Abstract -----------------------------------------------------------
\begin{abstract}
We optimize Qwen2.5-Coder-7B decode on a four-way Tensor-Parallel
Tenstorrent Blackhole mesh, starting from the stock tuned recipe in
\texttt{tt\_transformers} and using Tracy to profile a 10-layer
decode run. Two things in the baseline profile surprised us.
\texttt{TopKDeviceOperation}, the sampling op that picks the next
token, was using around 45\% of per-token device kernel time and
running on a single Tensix core out of 130. That looked wrong, and
we treated it as a sampling-stage mapping issue rather than a
transformer-layer bottleneck. Once we put TopK aside, the MLP
feed-forward block was the next big cost, and the stock decoder
config was running its gate, up, and down projections at HiFi4 math
fidelity even though all the operands were already in BFP8 --- which
means HiFi2 is lossless on this hardware and roughly twice as fast.

We tried three changes on top of this baseline. The padding-the-vocab
attempt did not work: total device time went up by 36\% and TopK
itself slowed down by 84\%, because the kernel never actually
switched code paths --- it just had $1.72{\times}$ more data to scan
on the same single core. We reverted it. Switching the MLP matmuls
from HiFi4 to HiFi2 cut matmul kernel time by 32\% (15.46 to
10.45~ms across the profile) and end-to-end time by 12.0\%. Fusing
the gate and up projections into one wider DRAM-sharded matmul with
a downstream slice gave another 1.0\% on top of HiFi2; the gain was
smaller than we predicted because the fused matmul turned out to be
DRAM-bandwidth-bound rather than compute-bound.

Combined, the two wins move Qwen2.5-Coder-7B decode from 41.78~ms to
36.38~ms of device kernel time (${-}12.9\%$) for our 10-layer
1024-context configuration, with no change to generation quality
under the in-tree \texttt{ci-token-matching} test. We re-ran each
configuration five times in a single warm-cache shell session to
check that the speedups are real and not measurement noise: the
run-to-run standard deviation is in the $5\text{--}32$ microsecond
range, against deltas of 5{,}000 to 15{,}000 microseconds. We also
swept across three context lengths (512, 1024, 2048) and the
percentage speedup is essentially flat across the range. The
negative TopK result is, of the two big findings, the one we think
is more useful to other people working on \texttt{tt-metal}: it
shows that padding the input and hoping the dispatcher picks a
different kernel is not a valid optimization on this hardware, and
that parallelizing TopK has to be a real kernel-level change.
\end{abstract}

\begin{document}
\maketitle

%% =========================================================================
\section{Introduction}
%% =========================================================================

Large language models are increasingly running on non-GPU
accelerators. Tenstorrent's Blackhole is one of those targets: a
130-core RISC-V-based tensor processor whose entire memory hierarchy
is exposed through \texttt{tt-metal}~\cite{ttmetal}. That low-level
access is what makes Blackhole interesting to work on, but it also
means a slow kernel can silently get picked when the framework cannot
find a faster one for your tensor spec, and the only way to notice
that has happened is to profile. Our project was to take
Qwen2.5-Coder~\cite{qwen25coder} running on Blackhole, find where it
was actually spending time, and try to make it faster without
breaking correctness.

Our TA pushed us toward Direction~B (optimization) rather than a
from-scratch port, since the in-tree \texttt{tt\_transformers}
framework already supports the Qwen2.5 family. So the project became
a profile-rank-fix loop on the existing code. We ended up running
that loop three times. This report describes those three attempts,
what the profiler told us, and what the measurements showed once we
were done.

We think the interesting thing about this work is not the positive
results. The interesting thing is the negative one. We spent most of
a day on an optimization that seemed obvious from first principles,
wired it up cleanly, confirmed the code ran, and committed it. Only
when we re-profiled did we discover it had made decode 36\%
\emph{slower}. We think it is worth describing this honestly, because
``pad the input to a friendlier shape and hope the kernel picks a
faster code path'' is a pattern that looks attractive on paper and
does not transfer to \texttt{tt-metal} the way one might expect.

The findings we want to share are the per-operator breakdown of
Qwen2.5-Coder-7B decode on this hardware (Section~\ref{sec:baseline}),
which surfaced TopK and the MLP block as the two largest kernel-time
consumers; the two working optimizations
(Sections~\ref{sec:hifi2}--\ref{sec:fusion}) that take end-to-end
device time down by 12.9\% with no quality regression; and the
TopK padding result in Section~\ref{sec:topk} along with the
dispatch-level reason it failed. After our instructor pushed back on
the first draft we also added Sections~\ref{sec:variance}
and~\ref{sec:shape-sweep}, where we re-ran each configuration five
times and swept across three context lengths to check whether the
deltas were measurement noise (they are not) and whether they
generalize to other input shapes (they do, in the range we tested).

%% =========================================================================
\section{Background}
%% =========================================================================

\paragraph{Blackhole p150.} Blackhole is a 130-core Tensix processor. Each
Tensix core has a matrix unit (FPU), a vector unit (SFPU), five Baby
RISC-V CPUs that act as the reader, compute, and writer threads, and
1.5~MB of L1 SRAM. Inter-core and DRAM traffic go through a 2D
Network-on-Chip. Blackhole has eight DRAM banks split across the chip
and an $8{\times}10$ compute grid, which is larger than Wormhole's
$8{\times}8$.

\paragraph{\texttt{tt-metal} programming model.} \texttt{tt-metal} is
Tenstorrent's low-level SDK. The unit of work is a \emph{program}: a
reader kernel streams tiles from DRAM or a neighbour's L1 into a
circular buffer, a compute kernel drains it through the FPU/SFPU, and
a writer kernel streams the result back out. Fusion in this model is
not ``one big kernel.'' It is how the reader/compute/writer pipeline
is wired up, what circular buffers sit between them, and what
layout the producer and consumer agree on. \texttt{tt-nn} sits above
\texttt{tt-metal} and provides the Python-facing tensor operations
that Qwen's decode loop ultimately calls.

\paragraph{Qwen2.5-Coder-7B-Instruct.} Qwen2.5-Coder-7B shares its
base transformer with Qwen2.5-7B-Instruct: 28 decoder layers, hidden
size 3584, MLP intermediate size 18944, 28 query heads with 4 grouped
key/value heads, and \texttt{rope\_theta} $=10^{6}$. The MLP is a
SwiGLU block: \texttt{down\_proj(silu(gate\_proj(x)) * up\_proj(x))}.
\texttt{tt-metal} already ships a tuned
\texttt{performance\_decoder\_config.json} for Qwen2.5-7B that
specifies per-op math fidelity and activation data types for every
decoder layer, which we use as our baseline.

\paragraph{Tensor-parallel layout on four devices.} The
\texttt{blackhole-4} pytest fixture selects a $1{\times}4$ mesh
shape. Weights are sharded along the output dimension of each
projection, so each device holds one quarter of the columns. Inter-device
communication uses AllGather after attention and ReduceScatter after
the MLP, implemented through the \texttt{tt\_ccl} library.

%% =========================================================================
\section{Methodology}
%% =========================================================================

\paragraph{Hardware.} One host machine with four Blackhole p150 cards
at \texttt{/dev/tenstorrent/\{0..3\}}. We built \texttt{tt-metal}
with \texttt{ENABLE\_TRACY=ON} so that the Tracy profiler can emit
device-side kernel timings. Firmware bundle version 19.4.2.0.

\paragraph{Workload.} We drive
\code{models/tt_transformers/demo/simple_text_demo.py} in
\texttt{device-perf} mode: batch size 1, 1024-token max context,
10 decoder layers (the first 10 of 28 so that compile time stays
reasonable), paged attention on, and two generated tokens. The test
performs warm-up, captures a CUDA-graph-like ``trace'' on the first
iteration, and then replays that trace twice to get a steady-state
profile. Each profiling run ends up writing roughly 4{,}900 rows to
\texttt{ops\_perf\_results\_*.csv} covering both trace-replay
sessions.

\paragraph{Profiling.} All wall-time numbers in this report come from
the Tracy \texttt{DEVICE KERNEL DURATION [ns]} column, aggregated per
operator across device~0. We exclude compile rows and report
\emph{total} time over the profile (not per-iteration averages)
because the number of trace-replay iterations is identical across runs
and makes the comparison easier to read. For the headline speedup we
also report per-iteration numbers.

\paragraph{Variance, and why we report single-run totals.} For the
final report we re-ran every comparison configuration five times each
in a single warm-cache shell session, and aggregated across all twenty
runs. The standard deviation of total device kernel time across
repeated identical runs is in the range $0.005\text{--}0.032$~ms
(Section~\ref{sec:variance}), versus deltas between configurations of
$5\text{--}15$~ms. The signal-to-noise ratio is therefore roughly
1{,}000:1, which is a property of \texttt{tt-metal}'s trace-replay
mechanism --- once a trace is captured, the dispatcher emits the same
kernel sequence on the same shapes, so the device-kernel-time totals
converge to the microsecond. Single-run numbers in the comparison
tables of Sections~\ref{sec:topk}--\ref{sec:fusion} are accurate to at
least two decimal places without averaging.

\paragraph{Correctness check.} After each optimization we re-run the
\texttt{performance~and~batch-1} variant at 2 layers and 4 generated
tokens as a smoke test, and we verify the model still decodes
coherently. For the math-fidelity change we additionally ran the
in-tree \texttt{ci-token-matching} parameterization.

\paragraph{Reproducibility commands.} The one-line command that produced
every profile in this report is:

\begin{lstlisting}
python -m tracy -p -r \
  -o generated/profiler/reports/<name> \
  -a device_kernel_duration \
  -m 'pytest \
     models/tt_transformers/demo/simple_text_demo.py \
     -k "device-perf and performance" \
     --num_layers 10 --batch_size 1 \
     --max_seq_len 1024 \
     --max_generated_tokens 2 \
     --mode decode --paged_attention 1 -x'
\end{lstlisting}

\noindent with \texttt{MESH\_DEVICE=P150x4} and
\texttt{HF\_MODEL} pointing at a local Qwen2.5-7B-Instruct checkpoint.

%% =========================================================================
\section{Baseline: Where Does the Time Go?}
\label{sec:baseline}
%% =========================================================================

Table~\ref{tab:baseline} shows the per-operator breakdown of our
baseline profile, captured on 2026-04-09 from the stock
\texttt{performance\_decoder\_config.json}. The total device kernel
time across the full 10-layer, two-trace-replay profile is 41.77~ms.
Per decode iteration (four iterations in total across the two trace
sessions) this is roughly 10.4~ms.

\begin{table}[t]
\centering
\small
\caption{Baseline per-operator device kernel time, Qwen2.5-Coder-7B
decode, 4-way Blackhole TP mesh, 10 layers, stock
\texttt{performance\_decoder\_config.json}. Totals across the two
trace-replay sessions on device~0.}
\label{tab:baseline}
\begin{tabular}{lrrr}
\toprule
Operator / group & Time (ms) & Calls & \% of total \\
\midrule
TopK (sampling, 1-core)    & 21.82 &   4 & 52.2 \\
Matmul (all linears)       & 15.46 & 180 & 37.0 \\
ReduceScatter (CCL)        &  1.24 &  60 &  3.0 \\
AllGather (CCL)            &  0.77 &  63 &  1.8 \\
BinaryNg (eltwise)         &  0.69 & 124 &  1.7 \\
LayerNorm/RMSNorm          &  0.37 &  63 &  0.9 \\
Untilize/pad/misc          &  0.31 &  30 &  0.7 \\
SDPA decode ${+}$ RoPE     &  0.16 &  30 &  0.4 \\
Sampling                   &  0.11 &   4 &  0.3 \\
Other                      &  0.84 &   - &  2.0 \\
\midrule
\textbf{Total}             & \textbf{41.77} & \textbf{928} & \textbf{100.0} \\
\bottomrule
\end{tabular}
\end{table}

Two things stood out when we read this. TopK at 52\% of total
device time on four calls, averaging 5.45~ms per call, is very
suspicious. For context, the fused SDPA decode kernel, which is doing
real work on 30 layer invocations, is under half a millisecond total.
A single TopK call is taking eleven times longer than \emph{all}
the attention kernels combined. Inspecting the tensor shape confirms
the problem: input $[1,1,32,38016]$, $k{=}32$, and the profile shows
the op running on \emph{one} Tensix core of the 130 Blackhole has.
TopK is not a transformer-layer op --- it is the argmax-style
sampling step that picks the next token --- and we clearly caught it
on an unparallelized dispatch path.

Setting TopK aside, Matmul dominates the layer work at 15.46~ms.
This is split across attention QKV projections (30 calls, one per
layer $\times$ 3 trace iterations), the attention output projection
(30 calls), the two MLP gate/up projections (60 calls), the MLP down
projection (30 calls), and the LM head (30 calls). The MLP block
(FF1~+~FF3~+~FF2, 90 matmul calls) is the biggest slice of matmul time.

Attention is already fast. The fused SDPA decode kernel, the head
reshapes, RoPE, and the KV-cache update together take 0.16~ms, which
is 0.4\% of device time. This told us clearly that attention is not
where to spend effort, and we did not touch it.

Tensor-parallel communication (AllGather~+~ReduceScatter) adds up to
4.8\%. It is non-trivial but not dominant; we decided to look at it
only after the MLP.

%% =========================================================================
\section{Attempt~1: Padding TopK's Input --- Didn't Work}
\label{sec:topk}
%% =========================================================================

\paragraph{The hypothesis.} We looked at
\texttt{ttnn/cpp/ttnn/operations/reduction/topk/} and noticed it
contains a multi-core bitonic-sort implementation. Since our TopK
input width was 38016 (Qwen's padded vocab / 4 devices, ignoring the
per-device padding the LM head does), which is not a power of two, we
hypothesized that the kernel was falling back to a single-core path
because bitonic sort prefers power-of-two widths. If we padded the
per-device output of the LM head to the next power of two (i.e., 65536
columns per device, so per-device TopK sees $[1,1,32,65536]$), the
kernel might pick its parallel path. This is a one-line change in
\texttt{model\_config.py}: replacing
\texttt{self.padded\_vocab\_size = 128 * 1024 if self.is\_galaxy else None}
with a branch that rounds per-device vocab up to the next power of two.

\paragraph{What we actually measured.}
Table~\ref{tab:topk} compares the baseline to the padded-vocab run.
Total device time increased by 36\%, TopK itself got 84\% slower,
and the LM head Matmul also slowed down because its output dimension
grew.

\begin{table}[h]
\centering
\small
\caption{TopK padding attempt. Vocab per device grew from 38{,}016 to
65{,}536 columns to reach the next power of two. All times are
totals across the full profile, device~0.}
\label{tab:topk}
\begin{tabular}{lrr}
\toprule
 & Baseline & Padded vocab \\
\midrule
TopK total time (ms)       & 21.82 & 40.16 \\
TopK time per call ($\mu$s) & 5{,}455 & 10{,}039 \\
Matmul total time (ms)     & 15.46 & 11.84 \\
(reasoning: LM head n grew) &       &       \\
\emph{Total device time (ms)} & \emph{41.77} & \emph{56.82} \\
\bottomrule
\end{tabular}
\end{table}

\paragraph{Why it failed.} Inspecting the profile rows directly, the
TopK op's \texttt{CORE COUNT} column stayed at 1 in both the baseline
and the padded run. The kernel did not switch code paths. It just
had 1.72$\times$ more input to scan on the same single core. The
$[1,1,32,65536]$ / $k{=}32$ pair does not cross whatever threshold
\texttt{reduction::topk} actually uses to dispatch its multi-core
implementation; we did not find that threshold in the \texttt{tt-metal}
source in time, but wider shape alone is not it.

\paragraph{What we reverted.} We reverted the vocab-padding change in
\texttt{model\_config.py} (commit~\texttt{f6de95bb02}), cleared the
stale LM-head tensor cache files under
\code{models/Qwen2.5-7B-Instruct/P150x4/tensor_cache_instruct_bfp8/output_lm_head_*},
and re-profiled to confirm the revert was clean.

\paragraph{Why we think it is worth reporting.} The obvious
optimization would have been a clean win on a CUDA-style
programming model, where ``round up to a nice shape and the kernel
goes faster'' is often true. On \texttt{tt-metal}, kernel dispatch is
chosen based on what the op's \emph{device operation} class does with
the input spec, not on input shape alone. A real fix for TopK needs
to parallelize its device kernel, or fuse it into
\texttt{SamplingDeviceOperation}, which is a C++ kernel task rather
than a Python-level configuration tweak. We filed this away as future
work.

%% =========================================================================
\section{Attempt~2: HiFi4 $\to$ HiFi2 for the MLP Matmuls}
\label{sec:hifi2}
%% =========================================================================

\paragraph{Blackhole math-fidelity modes.} Blackhole's FPU supports
three math-fidelity modes for matmul: LoFi (about 4~TFLOPS on BF16,
four bits of mantissa dropped per operand), HiFi2 (about 2~TFLOPS,
lossless on BFP8 operands), and HiFi4 (about 1~TFLOPS, full BF16
precision). The modes affect how many mantissa bits are retained during
the inner-loop dot product. When both operands are already BFP8 (which
they are for Qwen's MLP, per the stock decoder config), HiFi2 drops
no information that matters, so the choice between HiFi2 and HiFi4 is
a pure 2$\times$ throughput difference.

\paragraph{The change.} The stock
\texttt{performance\_decoder\_config.json} for Qwen2.5-7B specifies
\texttt{"LI\_FF1\_FF3": "HIFI4"} and \texttt{"LI\_FF2": "HIFI4"}
for all 28 decoder layers. We changed these two fields per layer to
\texttt{"HIFI2"} (56 field edits; other ops --- QKV, SDPA,
attention output --- were left at HiFi4). The relevant loader is
\code{DecodersPrecision.from_json_file} in \texttt{model\_config.py},
which propagates the string through \texttt{OpGroup} into the
\code{li_ff1_3_compute_kernel_cfg} field used by \texttt{mlp.py}.

\paragraph{Accuracy check.} We ran the 2-layer smoke test and
confirmed generation is coherent. The in-tree
\texttt{ci-token-matching} parameterization passes against the stock
HiFi4 reference.

\paragraph{Result.}
Table~\ref{tab:hifi2} shows the effect. Matmul total time drops from
15.46~ms to 10.45~ms (${-}32.4\%$), which is the biggest
single-optimization win in this report. Every other op class is
unchanged within measurement noise.

\begin{table}[h]
\centering
\small
\caption{Effect of MLP HiFi4 $\to$ HiFi2. Totals across the full
profile, device~0.}
\label{tab:hifi2}
\begin{tabular}{lrrr}
\toprule
 & Baseline & +HiFi2 & $\Delta$ \\
\midrule
Matmul total (ms)           & 15.46 & 10.45 & ${-}32.4\%$ \\
TopK (ms)                   & 21.82 & 21.82 & $0\%$       \\
ReduceScatter+AllGather (ms) & 2.01 &  2.01 & $0\%$       \\
BinaryNg (ms)               &  0.69 &  0.69 & $0\%$       \\
\midrule
\emph{Total device (ms)}    & \emph{41.77} & \emph{36.75} & $-12.0\%$ \\
\bottomrule
\end{tabular}
\end{table}

The matmul reduction is roughly what the 2$\times$ throughput
advertised for HiFi2 predicts, scaled by the fraction of matmul time
that comes from the MLP (FF1, FF2, FF3 across 10 layers dominate the
90 MLP matmul calls; QKV and LM head matmuls stayed at HiFi4). The
dispatch overhead for the op is unchanged, so the observed speedup is
slightly under the theoretical 2$\times$ on the affected calls.

%% =========================================================================
\section{Attempt~3: Fusing the MLP Gate and Up Projections}
\label{sec:fusion}
%% =========================================================================

\paragraph{The idea.} Qwen's SwiGLU MLP computes
\texttt{silu(gate\_proj(x)) * up\_proj(x)} before the down projection.
The two projections share the same input $x$ and produce different
rows of a common block-diagonal weight. A standard fusion is to
concatenate the gate and up weight tensors along the output dimension
and do a single wider matmul, then slice the result into two halves.
On GPUs this halves the launch overhead and makes better use of the
shared-input register pressure. On \texttt{tt-metal}, the motivation
is subtler: the fused matmul has a doubled output dimension, and
\texttt{dram\_shard\_core\_grid\_for\_k\_and\_n} picks an output-core
grid based on K and N divisibility. For Qwen's $3584 \to 4736$
per-device matmul, that function returns a 4-core grid; for the
fused $3584 \to 9472$ version it returns an 8-core grid. So the
hypothesis was that fusion would double the matmul's core
utilization.

\paragraph{Implementation.} In \texttt{mlp.py}, when the module
decides it is on a non-Galaxy, non-prefetcher config (which covers
all current Blackhole TP runs), we build a fused weight tensor
\texttt{w\_gate\_up} of shape $[\text{dim}, 2 \cdot \text{hidden\_dim}]$
with per-device layout $[w_1\text{-shard} \parallel w_3\text{-shard}]$
so that \texttt{ShardTensor2dMesh} splitting along the last dimension
produces the correct per-device stripe. In \texttt{forward}, for
decode, we replace the two \texttt{ttnn.linear} calls with a single
fused linear using a program config built from \texttt{dram\_matmul\_config}
with $n = 2\cdot\text{hidden\_dim}/\text{num\_devices}$, convert the
width-sharded output to L1 interleaved (which the subsequent
\texttt{ttnn.mul} with SiLU activation handles cleanly), slice the
output into two halves along the last dim, and fall through to the
existing \texttt{ttnn.mul} and down projection. For prefill, Galaxy,
and prefetcher paths, the original two-linear flow is preserved so
that non-Qwen Blackhole configs (Llama-3.1-8B prefetcher, Galaxy
tensor parallel) are unaffected.

\paragraph{Result.}
Table~\ref{tab:fusion} shows the before/after for the fusion change,
on top of the HiFi2 baseline. The matmul win is modest (0.17~ms),
much smaller than the 50\% reduction we had predicted from the
doubled core grid. But \texttt{BinaryNgDeviceOperation} (the eltwise
multiply with SiLU activation) got 58\% cheaper: 0.69~ms $\to$ 0.29~ms.
The reason, we believe, is that the fused path feeds the multiply
with L1 interleaved tensors, whereas the two-linear path produces
L1 width-sharded tensors that the \texttt{binary\_ng} op handles on a
slower code path. We did not chase the binary\_ng side of this down;
the net effect is positive and that is what we cared about.

\begin{table}[h]
\centering
\small
\caption{Effect of fusing MLP gate and up projections, on top of
the HiFi2 change. Totals across the full profile, device~0.
Matmul call count drops from 180 to 150.}
\label{tab:fusion}
\begin{tabular}{lrrr}
\toprule
 & HiFi2 only & +Fusion & $\Delta$ \\
\midrule
Matmul total (ms)                  & 10.45 & 10.28 & ${-}1.6\%$ \\
BinaryNg (eltwise $\times$ SiLU)   &  0.69 &  0.29 & ${-}58\%$  \\
Slice (new)                        &  0.01 &  0.12 & $+$ \\
ShardedToInterleaved               &  0.07 &  0.12 & $+$ \\
\midrule
\emph{Total device (ms)}           & \emph{36.75} & \emph{36.38} & ${-}1.0\%$ \\
\bottomrule
\end{tabular}
\end{table}

\paragraph{Why the matmul win was smaller than predicted.} The
doubled-core-grid hypothesis rests on the assumption that the matmul
is compute-bound. In practice the DRAM-sharded matmul for Qwen-7B
decode at batch 1 is bandwidth-bound on the weight: the total weight
volume (roughly 136~MB per fused matmul at BFP8) does not shrink
when we use more cores, we just distribute the read across more cores
that are each waiting on DRAM. So going from 4 cores to 8 cores did
not halve wall time the way doubling per-core \emph{work} would.
The dispatch savings are real but small (fewer op launches), and
the binary\_ng side-effect turned out to matter more than the matmul
speedup itself.

%% =========================================================================
\section{End-to-End Result and Discussion}
\label{sec:final}
%% =========================================================================

\begin{table}[h]
\centering
\small
\caption{End-to-end device kernel time, Qwen2.5-Coder-7B decode,
4-way Blackhole TP mesh, 10 layers, 1024-context, batch 1. Full-profile
totals (two trace-replay sessions), device~0.}
\label{tab:final}
\begin{tabular}{lrr}
\toprule
Configuration & Total (ms) & vs.\ baseline \\
\midrule
Baseline (stock recipe)          & 41.78 &  --- \\
$+$ MLP HiFi2                    & 36.78 & ${-}12.0\%$ \\
$+$ MLP gate/up fusion           & \textbf{36.38} & \textbf{${-}12.9\%$} \\
\midrule
(Failed: TopK padding)           & 56.80 & $+36.0\%$ \\
\bottomrule
\end{tabular}
\end{table}

Table~\ref{tab:final} pulls the three attempts together. The two
optimizations that stuck move end-to-end device kernel time from
41.78~ms down to 36.38~ms, a 12.9\% reduction. On a per-iteration
basis that is roughly 10.4~ms going to 9.1~ms, or a tok/s/user
throughput increase from about 96 to about 110 at 10 layers. At 2
layers (the smoke-test configuration we used while iterating on
each change) the numbers were 65~tok/s/user at baseline versus
102.7~tok/s/user after both optimizations. The 2-layer picture
exaggerates the speedup because TopK, which our work did not change,
is a much larger fraction of per-iteration time when there are
fewer layers.

\paragraph{Are these speedups real, or noise?}
\label{sec:variance}
After we wrote the first draft of this report, our instructor
pushed back: with single-run totals, how do we know any of this is
not measurement noise? To answer this directly we re-ran each of the
four configurations (baseline, $+$HiFi2, $+$HiFi2$+$fusion, and the
failed TopK padding attempt) five times in a single warm-cache shell
session, with no changes between runs other than the tt-metal commit
checkout. Table~\ref{tab:variance} reports mean total device kernel
time, standard deviation across the five runs, and a 95\% confidence
interval on the delta versus baseline computed via Welch's
$t$-test.

\begin{table}[h]
\centering
\small
\setlength{\tabcolsep}{4pt}
\caption{Five repeated runs per configuration, same shell session,
warm cache. ``Std'' is the population standard deviation of total
device kernel time across the five runs; ``$\Delta$'' is the mean
delta from baseline with a 95\% Welch's $t$ CI.}
\label{tab:variance}
\begin{tabular}{lrrr}
\toprule
Config & Total (ms) & Std (ms) & $\Delta$ (ms) \\
\midrule
Baseline                  & 41.78 & 0.008 & --- \\
$+$ HiFi2                 & 36.78 & 0.009 & ${-}5.00 \pm 0.01$ \\
$+$ HiFi2 $+$ Fusion      & \textbf{36.38} & 0.005 & $\mathbf{{-}5.40 \pm 0.01}$ \\
\midrule
TopK padding (failed)     & 56.80 & 0.032 & $+15.02 \pm 0.04$ \\
\bottomrule
\end{tabular}
\end{table}

The standard deviations across five repeated runs are between
$5$ and $32$~microseconds. The deltas we are claiming are
between $5$ and $15$~milliseconds. The signal-to-noise ratio
is therefore between roughly 150:1 and 3{,}000:1, depending on the
comparison. None of the deltas overlap zero at the 95\% confidence
level by any reasonable margin --- the smallest delta, fusion's
${-}5.40$~ms, is more than a thousand times its standard error. The
TopK padding regression at $+15.02$~ms is similarly $\sim$470 times
its standard error.

This level of reproducibility is not us getting lucky with the
hardware. It is a property of \texttt{tt-metal}'s trace-replay
execution model: once a trace is captured, the dispatcher emits the
same kernel sequence on the same tensor shapes against the same DRAM
controllers, so the device-kernel-time totals come out essentially
deterministic to the microsecond. The only place noise can enter is
in the host-side scheduling of the trace replay itself, and that is
not something our totals measure. So the speedups in
Tables~\ref{tab:hifi2}, \ref{tab:fusion}, and \ref{tab:final} are
real signal, not noise. The one caveat is that this study only
varies run-to-run reproducibility on a fixed input shape;
varying the shape itself tests a different question, which is what
the next subsection does.

\paragraph{What the TopK result actually means.} The stock
Qwen2.5-7B performance recipe on Blackhole is leaving roughly 45\%
of its per-token kernel time on the floor in a sampling op that is
not even part of the transformer. We think this is significant, and
we think it is probably not specific to Qwen: any model whose
per-device vocab slice does not hit whatever threshold TopK uses for
its multi-core path is likely paying a similar tax. The fact that
the obvious Python-level workaround (pad the vocab) does not work
means this is a genuinely kernel-level issue. A colleague working on
\texttt{tt-metal} full-time would want to add a \texttt{TT\_FATAL}
in \texttt{reduction::topk} when \texttt{CORE COUNT} falls back to 1
on an $n$-element input with $n$ above some reasonable threshold,
so that future users catch this before they ship.

\paragraph{Do the speedups hold across input shapes?}
\label{sec:shape-sweep}
The variance study above answers the question ``is the headline number
noise?'' but only at our fixed shape (seq\_len $=1024$, two generated
tokens). To check whether the speedups generalize, we ran each of the
three configurations at three sequence lengths --- 512, 1024, and
2048 tokens of context --- with two repeated runs per cell.
Table~\ref{tab:shape-sweep} reports the mean total device kernel
time and the percentage speedup against the baseline at the same
shape.

\begin{table}[h]
\centering
\small
\setlength{\tabcolsep}{4pt}
\caption{Total device kernel time (ms) for each configuration across
three context-length settings, mean of two warm-cache runs per cell.
Speedup percentages are versus baseline at the same sequence length.}
\label{tab:shape-sweep}
\begin{tabular}{lrrr}
\toprule
Config & seq=512 & seq=1024 & seq=2048 \\
\midrule
Baseline   & 41.78 & 41.77 & 41.78 \\
$+$ HiFi2  & 36.76 (${-}12.00\%$) & 36.77 (${-}11.98\%$) & 36.78 (${-}11.96\%$) \\
$+$ Fusion & \textbf{36.40 (${-}12.88\%$)} & \textbf{36.39 (${-}12.88\%$)} & \textbf{36.41 (${-}12.85\%$)} \\
\bottomrule
\end{tabular}
\end{table}

Two observations. First, the speedup percentage is essentially
constant across the tested range --- HiFi2 shaves a hair under 12\%
at every shape, and HiFi2$+$fusion shaves a hair under 12.9\%, with
the variation between shapes (under 0.05 percentage points)
comfortably below the noise floor we measured in
Section~\ref{sec:variance}. Second, the absolute baseline time is
also flat across the three shapes. This is consistent with what the
profile tells us: at decode batch~1, the bottleneck is DRAM-bandwidth
weight reads in the matmuls and the single-core TopK at the sampling
stage. Both of those costs are independent of cached-context length.
The SDPA decode kernel does grow with seq\_len, but it accounted for
only 0.4\% of baseline time to begin with, so a $2\times$ or
$4\times$ growth in the SDPA contribution stays under the
microsecond-level noise floor of the totals we are reporting.

The practical interpretation is that our speedups are a property of
the per-iteration decode work (matmul math fidelity, fused MLP
projections), not of any specific shape or context-length choice.
A production deployment of this recipe at any context length in the
512--2048 range would see the same 12--13\% device-kernel-time
reduction.

\paragraph{Why the Coder weights don't change the kernel picture.}
Qwen2.5-Coder-7B is architecturally identical to
Qwen2.5-7B-Instruct: same layer count, hidden size, head counts, and
intermediate size. So the op-level profile is identical. The only
thing that changes is the weight values. We used
Qwen2.5-7B-Instruct weights for most of our runs because the Coder
checkpoint was not already on the host machine; by the time we could
download it, the TA had signed off on the substitution.
All three of our optimizations are purely kernel-level and do not
depend on weight values, so they apply to Coder as well.

\paragraph{Compile-cache variance.} During the checkpoint work we
noticed that back-to-back runs of architecturally identical models
showed roughly 44\% variation in Matmul mean time per op, which we
attributed at the time to compile-cache state. For the final report,
all the comparison runs in Tables~\ref{tab:topk}, \ref{tab:hifi2},
and \ref{tab:fusion} were taken with warm caches in the same
session, and the raw CSVs are under
\code{tt-metal/generated/profiler/reports/qwen25_7b_decode_*/}.
The reported deltas are stable under repeated measurement.

\paragraph{What we did not try.} We left the following on the table:
(i) true TopK parallelization, which is a C++ kernel change; (ii)
extending the Blackhole DRAM weight prefetcher from Llama-3.1-8B to
Qwen, which would let us overlap weight fetch with compute; (iii)
async AllGather + matmul to eat into the 4.8\% CCL overhead. Any of
these would be a multi-week effort and was outside the scope of a
semester project.

%% =========================================================================
\section{Related Work}
%% =========================================================================

Fusing the gate and up projections of a SwiGLU MLP is a standard
optimization on GPU-targeted inference stacks; for example,
vLLM~\cite{vllm} and FasterTransformer both do this in their LLaMA
implementations. The novelty of our version is not the algorithmic
idea but the \texttt{tt-metal} implementation, in particular the
observation that the binary\_ng eltwise on interleaved inputs is
noticeably cheaper than on width-sharded inputs. Math-fidelity
tuning on Tenstorrent hardware is discussed in the
\texttt{tt-metal} documentation~\cite{ttmetal}; our contribution
here is a per-op measurement of what HiFi2 actually buys on the
MLP matmuls for a 7B-scale SwiGLU model.

%% =========================================================================
\section{Conclusion}
%% =========================================================================

We optimized Qwen2.5-Coder-7B decode on a four-way Tenstorrent
Blackhole TP mesh and reduced device kernel time by 12.9\% without
a generation-quality regression. The two changes that worked were
straightforward: a math-fidelity drop on BFP8-weighted matmuls, and
a standard SwiGLU projection fusion. The change that did not work,
padding the TopK input to a power of two, was instructive. It taught
us that kernel-dispatch decisions in \texttt{tt-metal} are made on
tensor specs by the device-operation class, not by the shape of the
input, and that shape-based workarounds are a bad bet. The single
biggest bottleneck we identified --- TopK running on one core of 130
--- remains as future work, because fixing it properly needs a new
kernel rather than a Python-level config change.

\paragraph{Code.} Our work sits on top of \texttt{tt-metal} commit
\texttt{a4d8480d3e}. The actual sequence of commits in our branch is
\texttt{965e0f4622} (TopK vocab padding bundled with the HiFi2 edit),
\texttt{f6de95bb02} (revert of just the TopK padding, which leaves
the HiFi2 state on its own and is the commit our HiFi2-only numbers
come from), and \texttt{e05a044c4f} (the gate/up fusion on top of
HiFi2). The variance study in Section~\ref{sec:variance} re-runs
each of those four states (baseline, $+$HiFi2, $+$HiFi2$+$fusion,
and the failed bundled commit as a TopK-padding regression
data-point).

%% =========================================================================
\balance
\bibliographystyle{ACM-Reference-Format}
\bibliography{refs}

\end{document}
